\scanpage{130}
\uline{NB} Le \inserted{MD-}graphe \deleted{\unsure{par}} $G'$ se déduit, à isom. unique près, de $G$ et de $C\inserted{=\vec C/\sigma}\subset A=\underbrace{A(G)}_{\substack{\text{ens. des arêtes}\\\text{liées}}}$.
Soit $\unsure{G_{0C}^{c}}$ le sous-graphe de $G$ défini par
\[
\begin{array}{l}
S(G_{0C})=S\\
\vec A(G_{0C})=\vec C\quad(\text{image inv. de }C\text{ dans }\widehat{\vec A})
\end{array}
\]
pondéré par $\uline{g}:\unsure{S}=S(G_{0C})\longrightarrow\mathbb N$ provenant de $G$.

On pose
\[
\left\{\begin{array}{l}
S'=\pi_0(G_{0C})\\
\unsure{\widehat{\vec A}'|\widehat A'}=\widehat{\vec A}\setminus\vec C\ \deleted{=\unsure{B}}\\
o':\widehat{\vec A}'\longrightarrow S'\quad\text{composé}\\
\sigma'\ \text{induit par }\sigma\\
g_{s'}=\text{genre virtuel de la comp. connexe}\\
\hfill G_{0C,s'}\ \text{du graphe pondéré }G_{0C}
\end{array}\right.
\]
\margin{où $\vec C$ est l'image inverse}
\diagram{\begin{tikzcd}[column sep=small,row sep=small]
\unsure{\widehat{\vec A}'}\arrow[r,hook]&\widehat{\vec A}\arrow[r,"o"]&S=S(G_{0C})\arrow[r,"\text{can.}"]&S'\arrow[d,equal]\\
&&&\pi_0(G_{0C})
\end{tikzcd}}

\uline{NB} On aura $I\simeq I'$, $\pi_0(G)\xrightarrow{\sim}\pi_0(G')$ et les applications $\gamma:\pi_0(G)-\mathbb N$ et $\gamma':\pi_0(G')-\mathbb N$ (genre virtuel d'une composante) se correspondent.

\mspage{129}
